\documentclass[11pt]{article}

\usepackage[spanish]{babel}
\usepackage[T1]{fontenc}
\usepackage[utf8]{inputenc}
\usepackage[margin=2.05cm]{geometry}
\usepackage{lmodern}
\usepackage{microtype}
\usepackage{amsmath,amssymb,mathtools}
\usepackage{enumitem}
\usepackage{booktabs}
\usepackage{xcolor}
\usepackage[hidelinks]{hyperref}

\definecolor{ink}{HTML}{1C2833}
\definecolor{rule}{HTML}{1F4E79}
\definecolor{boxbg}{HTML}{F4F7FA}
\definecolor{resultbg}{HTML}{EAF2F8}

\pagestyle{plain}
\setlength{\parskip}{0.45em}
\setlength{\parindent}{0pt}
\setlist[enumerate]{leftmargin=1.35em,itemsep=0.28em,topsep=0.3em}
\setlist[itemize]{leftmargin=1.15em,itemsep=0.22em,topsep=0.25em}

\newcommand{\curl}{\nabla\times}
\newcommand{\nv}{\nu}
\newcommand{\eg}{\varepsilon_{g}}
\newcommand{\bd}{\partial\Omega}
\newcommand{\dV}{\,\mathrm{d}\Omega}
\newcommand{\dS}{\,\mathrm{d}S}

\newcommand{\eqbox}[2]{%
  \par\smallskip
  \noindent\colorbox{resultbg}{%
    \begin{minipage}{\dimexpr\linewidth-2\fboxsep\relax}
      \vspace{0.2em}
      \[ #1 \tag{#2} \]
      \vspace{-0.45em}
    \end{minipage}}%
  \par\smallskip
}

\begin{document}

\begin{center}
{\Large\bfseries De Maxwell a la forma débil}\\[0.35em]
{\large Magnetostática y dinámica del banco de Kuntur}\\[0.55em]
{\small El camino que llega a \texttt{a(u,v)} y \texttt{l(v)}.
La \(\varepsilon\) del código no es la permitividad.}
\end{center}

\vspace{0.2em}
\noindent\textcolor{rule}{\rule{\linewidth}{0.6pt}}

\section*{1. Maxwell, y el corte que hace este problema}

En un medio lineal,
\begin{align}
\curl\mathbf{E} &= -\frac{\partial\mathbf{B}}{\partial t}, &
\curl\mathbf{H} &= \mathbf{J} + \frac{\partial\mathbf{D}}{\partial t}, \label{eq:maxwell}\\
\nabla\cdot\mathbf{B} &= 0, &
\nabla\cdot\mathbf{D} &= \rho, \nonumber
\end{align}
con \(\mathbf{B}=\mu\mathbf{H}\), \(\mathbf{D}=\varepsilon_0\mathbf{E}\) y
\(\mathbf{J}=\mathbf{J}_{s}+\sigma\mathbf{E}\).
La corriente \(\mathbf{J}_{s}\) es la que imponemos en la bobina.
La corriente \(\sigma\mathbf{E}\) es la que el campo induce en un conductor.

Kuntur modela aire, sin conductores en la malla, a frecuencias de hertz y en un dominio de decímetros. En ese régimen se apagan dos términos:

\begin{itemize}
  \item la corriente de desplazamiento \(\partial\mathbf{D}/\partial t\), porque el banco es diminuto frente a la longitud de onda (\(c/f\sim 3\times 10^{8}\,\mathrm{m}\) a \(1\,\mathrm{Hz}\));
  \item la corriente inducida \(\sigma\mathbf{E}\), porque en el aire \(\sigma=0\).
\end{itemize}

Lo que queda de \eqref{eq:maxwell} es la magnetostática
\begin{equation}
\curl\mathbf{H}=\mathbf{J}_{s},
\qquad
\nabla\cdot\mathbf{B}=0,
\qquad
\mathbf{H}=\nv\mathbf{B},
\qquad
\nv=\frac{1}{\mu_0}.
\label{eq:magnetostatica}
\end{equation}
El tiempo puede seguir estando en \(\mathbf{J}_{s}\). Si el operador no depende de \(t\), el campo lo sigue al instante. Esa es la dinámica que corre hoy el código.

\section*{2. Potencial y forma fuerte}

De \(\nabla\cdot\mathbf{B}=0\) existe un potencial vectorial tal que
\begin{equation}
\mathbf{B}=\curl\mathbf{A}.
\label{eq:B}
\end{equation}
Sustituir \eqref{eq:B} en \eqref{eq:magnetostatica} da la forma fuerte
\begin{equation}
\curl\bigl(\nv\,\curl\mathbf{A}\bigr)=\mathbf{J}_{s}
\quad\text{en }\Omega.
\label{eq:fuerte-singular}
\end{equation}

Esta ecuación no fija \(\mathbf{A}\). Si \(\mathbf{A}\) es solución, también lo es \(\mathbf{A}+\nabla\psi\), porque \(\curl\nabla\psi=\mathbf{0}\). El operador \(\curl\nv\curl\) tiene núcleo: los gradientes. La matriz de elementos finitos, sin un término extra, es singular.

El código corta ese núcleo con una penalización pequeña,
\begin{equation}
\boxed{\;
\curl\bigl(\nv\,\curl\mathbf{A}\bigr)+\eg\mathbf{A}=\mathbf{J}_{s}
\;}
\label{eq:fuerte}
\end{equation}
y en la frontera del aire
\begin{equation}
\mathbf{n}\times\mathbf{A}=\mathbf{0}
\quad\text{sobre }\bd.
\label{eq:bc}
\end{equation}
En el código, \(\eg\) es la constante \(\varepsilon=\nv\times 10^{-6}\). No es \(\varepsilon_0\). No impone \(\nabla\cdot\mathbf{A}=0\) de forma exacta: hace única la solución y empuja la parte de gradiente hacia cero. Si \(\eg\) es demasiado chica, \(K\) queda casi singular. Si es demasiado grande, contamina \(\mathbf{B}=\curl\mathbf{A}\).

\(\nv\) va dentro del rotacional, no fuera. En aire es constante y da igual; escrito así, la misma forma débil sirve el día que \(\nv\) cambie de un material a otro.

\section*{3. El paso a la forma débil}

El espacio donde vive la formulación es
\[
V=\bigl\{\mathbf{v}\in H(\mathrm{curl},\Omega):\ \mathbf{n}\times\mathbf{v}=\mathbf{0}\text{ en }\bd\bigr\},
\]
con \(H(\mathrm{curl})=\{\mathbf{v}\in L^{2}:\curl\mathbf{v}\in L^{2}\}\).
Ahí \(\curl\mathbf{v}\) existe en \(L^{2}\); no hace falta que \(\mathbf{v}\) tenga segunda derivada. Nédélec de orden 1 es un subespacio conforme de \(V\): la continuidad que impone es la tangencial, que es la que el rotacional necesita.

Se multiplica \eqref{eq:fuerte} por \(\mathbf{v}\in V\) y se integra:
\begin{equation}
\int_{\Omega}\mathbf{v}\cdot\curl(\nv\,\curl\mathbf{A})\dV
+\int_{\Omega}\eg\,\mathbf{v}\cdot\mathbf{A}\dV
=\int_{\Omega}\mathbf{v}\cdot\mathbf{J}_{s}\dV.
\label{eq:residual}
\end{equation}

La identidad que baja el rotacional, con \(\mathbf{w}=\nv\,\curl\mathbf{A}\), sale de
\(\nabla\cdot(\mathbf{v}\times\mathbf{w})=\mathbf{w}\cdot(\curl\mathbf{v})-\mathbf{v}\cdot(\curl\mathbf{w})\):
\begin{equation}
\int_{\Omega}\mathbf{v}\cdot(\curl\mathbf{w})\dV
=\int_{\Omega}(\curl\mathbf{v})\cdot\mathbf{w}\dV
-\int_{\bd}(\mathbf{v}\times\mathbf{w})\cdot\mathbf{n}\dS.
\label{eq:partes}
\end{equation}
El borde se anula porque \(\mathbf{n}\times\mathbf{v}=\mathbf{0}\). El integrando cumple
\[
(\mathbf{v}\times\mathbf{w})\cdot\mathbf{n}=(\mathbf{n}\times\mathbf{v})\cdot\mathbf{w}=0.
\]
\(\mathbf{v}\times\mathbf{w}\) no tiene por qué ser cero: queda tangente a \(\bd\), y lo que entra en la integral es su componente normal. No hay integral de frontera que ensamblar. La condición \eqref{eq:bc} ya está metida en el espacio, que es el \texttt{dirichlet\_tags = ["boundary"]} de Nédélec.

Sustituir \(\mathbf{w}=\nv\,\curl\mathbf{A}\) en \eqref{eq:residual} deja la forma débil. Hallar \(\mathbf{A}\in V\) tal que, para todo \(\mathbf{v}\in V\),
\eqbox{
\int_{\Omega}\nv\,(\curl\mathbf{A})\cdot(\curl\mathbf{v})\dV
+\int_{\Omega}\eg\,\mathbf{A}\cdot\mathbf{v}\dV
=\int_{\Omega}\mathbf{v}\cdot\mathbf{J}_{s}\dV.
}{\(*\)}
Eso es, línea por línea, lo que declaran los tres planteamientos
(\texttt{helmholtz.jl}, \texttt{ring.jl}, \texttt{triaxial.jl}):
\begin{equation*}
\begin{aligned}
a(u,v)&=\int\bigl(\nv\,(\nabla\times u)\cdot(\nabla\times v)+\varepsilon\,(u\cdot v)\bigr)\,\mathrm{d}\Omega,\\
l(v)&=\int v\cdot J\,\mathrm{d}\Omega.
\end{aligned}
\end{equation*}
El primer término es la energía magnética. El segundo es el gauge. El lado derecho es el trabajo de la corriente impuesta.

\(a(u,v)=a(v,u)\): la matriz es simétrica. Con \(\eg>0\) es definida positiva. Sin \(\eg\) solo es semidefinida, y el núcleo discreto (gradientes de Nédélec) deja el sistema sin solución única.

\section*{4. Estática: un sistema lineal}

La corriente no depende del tiempo. \(\mathbf{J}_{s}\) es el par de anillos (o los tres pares) a \(I_{0}\) amperes, con la gaussiana azimutal normalizada para que la integral en la sección valga esa corriente.

Nédélec escribe \(\mathbf{A}_{h}=\sum_{j} a_{j}\boldsymbol{\varphi}_{j}\).
Probar contra cada \(\boldsymbol{\varphi}_{i}\) convierte \((\ast)\) en
\begin{equation}
Ka=f,
\label{eq:sistema}
\end{equation}
\begin{align}
K_{ij}
&=a(\boldsymbol{\varphi}_{j},\boldsymbol{\varphi}_{i})
=\int_{\Omega}\nv\,(\curl\boldsymbol{\varphi}_{j})\cdot(\curl\boldsymbol{\varphi}_{i})\dV
+\int_{\Omega}\eg\,\boldsymbol{\varphi}_{j}\cdot\boldsymbol{\varphi}_{i}\dV,
\label{eq:K}\\
f_{i}
&=l(\boldsymbol{\varphi}_{i})
=\int_{\Omega}\boldsymbol{\varphi}_{i}\cdot\mathbf{J}_{s}\dV.
\label{eq:f}
\end{align}
\texttt{AffineFEOperator(a, l, U, V)} ensambla \eqref{eq:sistema} y lo resuelve una vez. Del coeficiente se recupera el campo; el rotacional es la derivada que el espacio ya sabe calcular:
\begin{equation}
\mathbf{B}_{h}=\curl\mathbf{A}_{h},
\qquad
\mathbf{H}_{h}=\frac{1}{\mu_{0}}\mathbf{B}_{h}=\nv\mathbf{B}_{h}.
\label{eq:BH}
\end{equation}
La energía que imprime el programa es la del término de rigidez, sin el gauge:
\begin{equation}
W=\frac12\int_{\Omega}\mathbf{B}_{h}\cdot\mathbf{H}_{h}\dV.
\label{eq:energia}
\end{equation}
Nada de esto se vuelve a factorizar si \(I_{0}\) cambia: como \(f\) es lineal en la corriente, basta escalar \(a\), \(\mathbf{B}_{h}\) y \(\mathbf{H}_{h}\).

\section*{5. Dinámica}

Hay tres dinámicas distintas. Solo la primera está en el código. Las otras dos cambian la forma débil, no solo el lado derecho.

\subsection*{5.1 Aire, operador fijo: escalar \(I(t)\)}

En aire \(\sigma=0\) y \(\nv\) no depende del tiempo. La corriente entra solo en \(f\):
\begin{equation}
\mathbf{J}_{s}(\mathbf{x},t)=I(t)\,\hat{\jmath}(\mathbf{x}),
\label{eq:Jt}
\end{equation}
con \(\hat{\jmath}\) la misma fuente ya normalizada a \(1\,\mathrm{A}\).
La forma débil en cada instante es \(a(\mathbf{A}(t),\mathbf{v})=I(t)\,l_{1}(\mathbf{v})\).
Como \(a\) es lineal e independiente de \(t\),
\begin{equation}
\boxed{\;
\mathbf{A}(\mathbf{x},t)=I(t)\,\mathbf{A}_{1}(\mathbf{x}),
\qquad
\mathbf{B}(\mathbf{x},t)=I(t)\,\mathbf{B}_{1}(\mathbf{x})
\;}
\label{eq:escala}
\end{equation}
donde \((\mathbf{A}_{1},\mathbf{B}_{1})\) es la solución estática de la sección 4 con \(I=1\,\mathrm{A}\).

Eso es \texttt{main\_transient.jl}: un solo \(Ka=f\), y cada cuadro del VTK multiplica el campo por \(I(t)/I_{0}\). El suavizado gaussiano de \(I(t)\) no entra en la forma débil; solo fabrica la escala. Una senoide, un escalón o una tabla producen el mismo \(\mathbf{B}_{1}\). Cambia el número por el que se multiplica.

Esto vale mientras el campo alcance el régimen de \eqref{eq:magnetostatica} mucho antes de que \(I(t)\) cambie. A \(1\,\mathrm{Hz}\), en aire y en \(10\,\mathrm{cm}\), se cumple de sobra: no hay onda dentro del dominio y no hay conductor donde pueda circular una corriente inducida.

\subsection*{5.2 Conductor: la forma débil depende de la historia}

Si el dominio contiene un conductor, Faraday ya no se puede ignorar. Con \(\mathbf{B}=\curl\mathbf{A}\) y el gauge \(\phi=0\),
\begin{equation}
\mathbf{E}=-\frac{\partial\mathbf{A}}{\partial t},
\qquad
\mathbf{J}=\mathbf{J}_{s}+\sigma\mathbf{E}.
\label{eq:E}
\end{equation}
Ampère, otra vez sin \(\partial\mathbf{D}/\partial t\), queda
\begin{equation}
\curl\bigl(\nv\,\curl\mathbf{A}\bigr)+\sigma\frac{\partial\mathbf{A}}{\partial t}=\mathbf{J}_{s}.
\label{eq:eddy-fuerte}
\end{equation}
El mismo paso de la sección 3, con el mismo borde \(\mathbf{n}\times\mathbf{A}=\mathbf{0}\), da
\eqbox{
\int_{\Omega}\nv\,(\curl\mathbf{A})\cdot(\curl\mathbf{v})\dV
+\int_{\Omega}\sigma\frac{\partial\mathbf{A}}{\partial t}\cdot\mathbf{v}\dV
=\int_{\Omega}\mathbf{v}\cdot\mathbf{J}_{s}\dV.
}{\(**\)}
El primer término es la misma \(K\) de \eqref{eq:K}. El segundo es nuevo: una masa solo donde \(\sigma>0\),
\begin{equation}
M_{ij}=\int_{\Omega}\sigma\,\boldsymbol{\varphi}_{j}\cdot\boldsymbol{\varphi}_{i}\dV.
\label{eq:M}
\end{equation}
En aire, \(\sigma=0\) y esa integral no aporta. El sistema semi-discreto es
\begin{equation}
M\frac{\mathrm{d}a}{\mathrm{d}t}+Ka=f(t).
\label{eq:ode}
\end{equation}
Ya no basta calcular \(\mathbf{A}_{1}\) y escalar. \(a(t)\) depende de \(a\) en instantes anteriores: una corriente que se apaga deja campo remanente mientras las corrientes de Foucault se disipan.

Euler implícito, \(\mathrm{d}a/\mathrm{d}t\approx(a^{n+1}-a^{n})/\Delta t\), convierte \eqref{eq:ode} en un sistema del mismo tipo que \eqref{eq:sistema}, reensamblado en el lado izquierdo:
\begin{equation}
\Bigl(\frac{M}{\Delta t}+K\Bigr)a^{n+1}
=\frac{M}{\Delta t}\,a^{n}+f^{n+1}.
\label{eq:euler}
\end{equation}
\(K\) se ensambla una vez. Cambia el lado derecho, y \(M/\Delta t\) se suma a \(K\).

\subsection*{5.3 Onda completa}

Si se restituye \(\partial\mathbf{D}/\partial t\) y se mantiene \(\mathbf{E}=-\partial\mathbf{A}/\partial t\), la forma fuerte gana un término de segundo orden,
\begin{equation}
\curl\bigl(\nv\,\curl\mathbf{A}\bigr)
+\sigma\frac{\partial\mathbf{A}}{\partial t}
+\varepsilon_{0}\frac{\partial^{2}\mathbf{A}}{\partial t^{2}}
=\mathbf{J}_{s},
\label{eq:onda}
\end{equation}
y la forma débil, el término
\(\int\varepsilon_{0}\,(\partial^{2}\mathbf{A}/\partial t^{2})\cdot\mathbf{v}\,\mathrm{d}\Omega\).
Esa \(\varepsilon_{0}\) sí es la permitividad. A \(1\,\mathrm{Hz}\) en este banco es irrelevante frente a \(K\). Entra cuando el tamaño del dominio deja de ser despreciable frente a \(c/f\).

\section*{6. Mapa}

\begin{center}
\renewcommand{\arraystretch}{1.35}
\begin{tabular}{@{}llll@{}}
\toprule
Régimen & Forma débil & Sistema & En Kuntur \\
\midrule
Estática & \((\ast)\) con \(\mathbf{J}_{s}\) fijo & \(Ka=f\) & \texttt{main.jl} \\
Aire, \(I(t)\) & \((\ast)\) con \(\mathbf{J}_{s}=I(t)\hat{\jmath}\) & \(a(t)=I(t)\,a_{1}\) & \texttt{main\_transient.jl} \\
Foucault & \((\ast\ast)\) & \(M a'+Ka=f(t)\) & no está \\
Onda & \((\ast\ast)\) más \(\varepsilon_{0}\partial_{tt}\) & \(M_{2}a''+Ma'+Ka=f\) & no está \\
\bottomrule
\end{tabular}
\end{center}

\end{document}
